%!TEX program = uplatex
\RequirePackage{plautopatch}
\documentclass[uplatex,dvipdfmx]{jlreq}
\usepackage{amsmath,amssymb}
\usepackage{enumerate}

\pagestyle{empty}

\begin{document}

%======================================================================
% 講演1：小野薫「運動量写像とシンプレクティック簡約」
%======================================================================
\begin{center}
{\large\bfseries 運動量写像とシンプレクティック簡約}\\[6pt]
\hspace*{\fill}{\large\bfseries 小野　薫（北大・理）}
\end{center}
\vspace{1pc}

この講演では、例を通してシンプレクティック簡約とは何かを
知って頂き、後の講演への warm-up となることを目標としたい。
シンプレクティック構造は、解析力学の記述などで知っている方々も
居られると思う。古典力学において、系の対称性と保存則が表裏一体である
ことは Noether の定理として知られている。講演ではこのあたりから始めて
シンプレクティック構造や運動量写像、また運動量写像を持つときに、
シンプレクティック簡約と呼ばれる構成を例を用いながら説明したい。
時間が許せば、Hermite 行列の固有値に関わる話題を題材に、シンプレクティック簡約
に関わる事項（運動量写像の凸性、安定性、Borel--Weil の定理など）の絡み合いに
ついても話したい。

\newpage

%======================================================================
% 講演2：森吉仁志「シンプレクティック多様体と量子化」
%======================================================================
\begin{center}
{\large\bfseries シンプレクティック多様体と量子化}\\[6pt]
\hspace*{\fill}{\large\bfseries 森吉　仁志（慶応大・理工）}
\end{center}
\vspace{1pc}

この講演では、参加者に「シンプレクティック多様体の量子化」を体験していただく
ことを目標としたい。シンプレクティック多様体においては、無限回微分可能な関数
全体のなす関数環（点毎の掛け算により積を定める）がポアソン括弧というリー環の
構造をもつ。量子化の観点からは、ポアソン括弧は次のように解釈できる。まず新た
な非可換積を関数環に与える。これは関数を作用素として再定義（量子化）し、作用
素の積を関数環に導入することで行われる。このとき量子化の操作にパラメター（プ
ランク定数）を組み込んでおく。ここで、パラメターが0に近づくとき（古典近
似）、非可換積は関数環の可換な積に回復されることを要請する。次に作用素の交換
子をパラメターで展開する。このときパラメターに関する一次近似として現れるリー
括弧がポアソン括弧になる。このような定式化によれば、ポアソン括弧とは、量子化
（非可換化）を行った後に関数環に残されるよすがと考えられる。

これは数ある量子化のひとつの定式化であるが、講演ではこの定式化に沿うさまざ
まな量子化の方法に関して、具体例を紐解きながら解説を行う。とくに、簡明な空間
であるユークリッド空間や双曲空間を取り上げて、正準交換関係を用いた量子化（ハ
イゼンベルグの量子化）、フォック空間による量子化、幾何的量子化、テープリッツ
作用素による量子化、Berezin 量子化、Moyal 積による量子化、Weyl 量子化などをお見
せするつもりである。

言うまでもないが、講演者は量子化あるいはシンプレクティック幾何学の専門家で
は到底ない。しかし非専門家としてシンプレクティック多様体に関わる様子を、皆様
の体験に供するようお話しできれば幸いである。

\newpage

%======================================================================
% 講演3：高倉樹「シンプレクティック商のトポロジー」
%======================================================================
\begin{center}
{\large\bfseries シンプレクティック商のトポロジー}\\[6pt]
\hspace*{\fill}{\large\bfseries 高倉　樹（中央大・理工）}
\end{center}
\vspace{1pc}

ハミルトニアンな群作用と運動量写像、および
シンプレクティック商の概念は、
対称性をもつハミルトン力学系に由来する長い歴史を持つ。
反面、設定がやや特殊であるせいもあり、
現在のシンプレクティック幾何学における位置付けとしては、
比較的マイナーな対象ともいえる。
とはいえ80年代以降の研究を振り返ってみると、
（トポロジーに関する話題に限るとしても）
それなりに大きな流れがある。序に代えて、
私見を少々述べよう。

80年代前半、当時
のシンプレクティック幾何学固有の問題意識を踏まえて、
幾何学的量子化に関する Guillemin--Sternberg の論文 [3] が現れる。
シンプレクティック商の重要性を明示する画期的な結果といえる。
一方ほぼ同じ頃、新たな波と呼ぶにふさわしい形で、
2次元ゲージ理論に関する
Atiyah--Bott の論文 [1] と、
その有限次元版である Kirwan の著作 [2] が登場する。
単にシンプレクティック商のトポロジーを調べる手法が得られた
だけではなく、後の研究につながる多くの示唆が与えられている。
なお、ここでは上記2つを独立な結果であるかのように述べたが、
運動量写像の凸性や代数群作用に関する安定性との対応などの
本質的な現象が、それぞれの立場から捉えられていることを注意しておく。

その後90年前後から、（位相的）場の量子論をさらなる動機として加えつつ、
シンプレクティック商のトポロジーの精密な考察は
2次元ゲージ理論とともに発展する。
同時に上記2つの流れはさらに交錯を深めていくことになる。
この点に関し、Witten の果たした役割は大きい（[4] などを参照）。
\vspace{1pc}

さて、そもそも本講演はこの分野のサーベイを意図するものではない。
トポロジーへの応用に関していうならば、これらの理論は、
特殊ではあるが興味深い諸例を念頭において築かれてきた。
ここでも（ゲージ理論からはひとまず離れつつ）典型的な例を用いて、
諸定理の一端を紹介する。

ハミルトニアンな群作用を持つ多様体の例としては、
定義により、リー群の余随伴軌道が基本的である。
（余随伴軌道とは聞きなれない用語かも知れぬ。多くの文脈では、
旗多様体と呼ばれる。）
またそれらから派生する空間は、不変式論や表現論と直接関わる。
なお、空間多角形のモジュライ空間は、このようにして得られる
シンプレクティック商の典型例であるが、
それは点つき球面上の平坦接続のモジュライ空間の特別な場合とも見なされる。
これらの例を手にとって、現象を実感することを目標としたい。
\vspace{1pc}

\begin{center}\textsc{参考文献}\end{center}
\begin{enumerate}[{[1]}]
\item
M.~F.~Atiyah and R.~Bott, \textit{The Yang--Mills equations over Riemann
surfaces}, Phil. Trans.~Roy.~Soc.~Lond. \textbf{A308} (1982) 523--615.
\item
F.~Kirwan, \textit{Cohomology of Quotients in Symplectic and Algebraic Geometry}, Princeton University Press (1984).
\item
V.~Guillemin and S.~Sternberg,
\textit{Geometric quantization and multiplicities of group representations},
Invent.~Math. \textbf{67} (1982) 515--538.
\item
E.~Witten, \textit{Two dimensional gauge theories revisited},
J.~Geom.~Phys. \textbf{9} (1992) 303--368.
\end{enumerate}

\newpage

%======================================================================
% 講演4：古田幹雄「曲率と運動量写像 I, II」
%======================================================================
\begin{center}
{\large\bfseries 曲率と運動量写像 I, II}\\[6pt]
\hspace*{\fill}{\large\bfseries 古田　幹雄（東大・数理）}
\end{center}
\vspace{1pc}

2回の話で、シンプレクティック幾何に現れる
運動量写像が、低次元幾何学において
接続の曲率として現れることを説明したいと思います。

無限次元の空間に有限次元の多様体に対する理論が
形式的に適用可能であり、非自明な結果が得られる
典型例となっています。ポイントは、一度無限次元の空間の
考察を経由しているにも関わらず、
シンプレクティック商として、有限次元の多様体（しばしば特異点をもつ）が得られ
ることです。
すると、この有限次元の多様体は自然なシンプレクティック
構造を持つことがわかります。
この場合、有限次元の多様体を性質を自然な枠組みで
理解するために無限次元の空間を
一度経由することが必要となります。

講演者はこの分野の専門家ではありませんが、1980年代に
発展した基本的な部分について、次の順序で説明する予定です。

\begin{itemize}
\item
接続と曲率の定義
\item
2次元のゲージ理論とモーメント写像
\item
4次元のゲージ理論とモーメント写像
\end{itemize}

2次元の向きのある閉多様体の上に
主束がひとつ与えられたとします。
このとき、この主束の接続全体の空間には、
シンプレクティック構造がはいります。
接続全体の空間には
主束の自己同型群（ゲージ群）が作用しており、
このシンプレクティック構造を保ちます。
このとき、各接続に対してその曲率を対応させる写像は、
モーメント写像として解釈することができます。
さらに底空間に複素構造、すなわち Riemann 面としての
構造を与えると、それに伴って接続全体の空間には
無限次元の Kähler 多様体の構造がはいり、
ゲージ群の複素化が、複素構造を保つ作用をもちます。

これが Atiyah--Bott による Riemann 面上の Yang--Mills 場の研究の
出発点でした [AB]。
モーメント写像の一般論をこの無限次元の状況に
いかに適用するかが問題であるといえます。
しかし歴史的には、
この無限次元の状況を理解する道具の開発を通じて、
幾人もの人々によって一般論が、
構築されてきたというのが筋道でした。

上の構成はいくつかの方向に拡張されていますが、
講演では、ひとつの方向として次のことを紹介する予定です：

4次元の向きのある閉多様体上に
主束がひとつ与えられたとします。
もしその4次元多様体が超 Kähler 構造をもつなら、
接続全体の空間も超 Kähler 構造をもち、
曲率の自己双対部分は、超 Kähler 構造に対する
モーメント写像として解釈することが可能となります。
このとき超 Kähler 商は ASD 接続のモジュライ空間となり、
その上に超 Kähler 構造の構造が自然に入ります。

\vspace{3mm}

\noindent
参考文献

\noindent
シンプレクティック商について

\noindent
[AB]
Atiyah, M.~F., Bott, R.,
``The Yang--Mills equations over Riemann surfaces''.
Philos. Trans. Roy. Soc. London Ser. A 308 (1983),
no. 1505, 523--615.

\noindent
[K]
Kirwan, F.,
``Cohomology of quotients in symplectic and algebraic
geometry'', Math. Notes 31, 1984,
Princeton Univ. Press
\vspace{6pt}

\noindent
超ケーラー商について

\noindent
[HKLR]
Hitchin, N.~J., Karlhede, A.,
Lindström, U., Roček, M.,
``Hyper-Kähler metrics and supersymmetry'',
Comm. Math. Phys. 108 (1987), no. 4, 535--589.

\noindent
[AH]
Atiyah, M.~F., Hitchin, N.~J.
``Low energy scattering of nonabelian monopoles'',
Phys. Lett. A 107 (1985), no. 1, 21--25.

\noindent
[H]
Hitchin, N.~J.
``The self-duality equations on a Riemann surface'',
Proc. London Math. Soc. (3) 55 (1987), no. 1, 59--126.

\noindent
[I]
Itoh, M.
``Quaternion structure on the moduli space of
Yang--Mills connections'',
Math. Ann. 276 (1987), no. 4, 581--593.

\newpage

%======================================================================
% 講演5：太田啓史「シンプレクティック商と安定性／ゲージ理論と安定性」
%======================================================================
\begin{center}
{\large\bfseries 1. シンプレクティック商と安定性\\
2. ゲージ理論と安定性}\\[6pt]
\hspace*{\fill}{\large\bfseries 太田　啓史（名大・多元数理）}
\end{center}
\vspace{1pc}

\noindent
\textbf{1.} コンパクト群によるシンプレクティック商は、
しばしば代数的な幾何学的不変式（Geometric Invariant Theory）商と
対応することがあります。
コンパクト群による商空間は、分離性については大して問題はありませんが、
不変式商の方は、いわばその複素化された非コンパクト群による商を
扱うことになり、その商空間の分離性は極めて基本的な問題となります。
この時、安定性が重要な役割を果たします。
安定性は、そのひとつの見方として、
商がしかるべき分離性をもつために現れると考えられます。
ここでは、
シンプレクティック商と幾何学的不変式商の対応において
安定性が果たす役割について、
有限次元の具体的な例を通してお話する予定でいます。
ひとつのものを2通りの立場で見ることの簡単な例になると思います。
\vspace{1pc}

\noindent
\textbf{2.}
1でお話した有限次元の場合のアナロジーを無限次元の場合に考えると、
ときに、非自明な結論を導くことがあります。
一般にはそのままでは成り立たないかもしれませんし、
成り立つとしてもその証明には、新たな道具の開発などが必要となり
決して安直なアナロジーで証明ができるものではなく、それこそが
重要な点であると強調したいところですが、残念ながら講演では踏み込みません。
ここでは、リーマン面上の接続
あるいは、ヒッグス場を組みにした状況での2つの商の関係を
考察します。それぞれの商の幾何学的構造を行き来することによって、
例えば、安定性の一つの帰結としてリーマン面上の平坦束に関する
Milnor--Wood の不等式などがひょっこり出てくることを紹介できればと思います。
時間があれば、安定性を基軸として
他のいくつかの不等式も俯瞰できればと思います。
\vspace{6pt}

\noindent
参考文献：

\noindent
Hitchin, N.~J.
``The self-duality equations on a Riemann surface'',
Proc. London Math. Soc. (3) 55 (1987), no. 1, 59--126.

\end{document}